\documentclass[11pt]{article}
\usepackage{hyperref}
\usepackage{enumitem}
\usepackage{cmap}
\usepackage[utf8]{inputenc}
\usepackage[T1]{fontenc}
\usepackage[left=1.15cm,top=1.35cm,right=1.15cm,bottom=0.9cm]{geometry}
\IfFileExists{glyphtounicode.tex}{\input{glyphtounicode}\pdfgentounicode=1}{}
\hypersetup{
    hidelinks,
    pdfauthor={Kevin Bravo},
    pdftitle={Kevin Bravo - Senior Backend Developer (Node.js, TypeScript, IA en Produccion)},
    pdfsubject={CV para el rol de Sr. Backend Developer en Retorna: backend Node.js y TypeScript de alta criticidad, IA en produccion real, agentes autonomos, arquitectura multi-tenant y ownership end to end},
    pdfkeywords={Kevin Bravo, Senior Backend Developer, Node.js, TypeScript, IA en produccion, agentes autonomos, LLM, guardrails, human-in-the-loop, PostgreSQL, multi-tenant, arquitectura distribuida, Go, Python, fintech}
}
\setlength{\parindent}{0pt}
\setlist[itemize]{leftmargin=*, noitemsep, topsep=2pt, partopsep=0pt, parsep=2pt, label={\textbf{-}}}

\begin{document}

\begin{center}
    {\LARGE \textbf{Kevin Bravo}}\\
    Senior Backend Developer | Node.js, TypeScript, PostgreSQL, AWS \\
    \hrulefill
\end{center}

\begin{center}
    \href{mailto:0bravokevin@gmail.com}{0bravokevin@gmail.com}
    \textbullet\
    \href{https://kevinbravo.com}{kevinbravo.com}
    \textbullet\
    \href{https://github.com/0bkevin}{github.com/0bkevin}
    \textbullet\
    \href{https://www.linkedin.com/in/0bkevin/}{linkedin.com/in/0bkevin}
\end{center}

\vspace{4pt}
\begin{center}
    \textbf{Perfil Profesional}
\end{center}
\vspace{-4pt}
Ingeniero backend con 5 años de experiencia construyendo y operando sistemas en Node.js, TypeScript, Python y Go. He sido responsable del ciclo completo, desde los requisitos y el modelo de datos hasta el despliegue, los releases y el soporte en producción. Como único ingeniero de Billingua Talent, construí una plataforma multi-tenant con acceso por roles y aislamiento de datos para cuatro tipos de usuario. En AVAA desarrollé SEP y SEA, sistemas que centralizan procesos operativos, pagos, recibos y reportes; SEP atiende a cerca de 300 becarios en tres capítulos nacionales. También lideré una migración de React a SvelteKit sobre un producto activo y desarrollo sistemas de agentes de IA con enrutamiento, control de sesiones, registros de ejecución y monitoreo. Me interesan los sistemas distribuidos donde la trazabilidad, la integridad de los datos y la corrección son esenciales.

\vspace{8pt}
\begin{center}
    \textbf{Habilidades Técnicas}
\end{center}
\vspace{-4pt}
\textbf{Backend:} Node.js, TypeScript, JavaScript, Python, Django, Go, REST APIs, diseño de APIs, arquitectura de servicios, integraciones con terceros, ownership end to end \\
\textbf{IA aplicada:} APIs de LLM (Anthropic, OpenAI), agentes autónomos, planos de control de agentes, tool use, routing, MCP, control de sesiones, logs de ejecución, health checks y supervisión humana \\
\textbf{Datos y arquitectura:} PostgreSQL, Prisma, SQL, modelado de datos relacional, arquitectura multi-tenant, control de acceso por rol, aislamiento por tenant, logging de acciones y observabilidad, ETL \\
\textbf{Cloud y entrega:} Azure, VPS, Cloudflare Workers, Docker, despliegues, debugging en producción, releases y entrega continua

\vspace{8pt}
\begin{center}
    \textbf{Sistemas de IA en Producción}
\end{center}
\vspace{-4pt}
\textbf{Brio} (Go, React Native, PostgreSQL): Plano de control para agentes de IA, construido como un sistema distribuido con un servicio companion, un relay y un protocolo compartido. Maneja emparejamiento de dispositivos, enrutamiento, propiedad por usuario, control de sesiones, logs y health checks. \\
\textbf{Zrode} (TypeScript): Mi propia versión de una GUI open-source para agentes de coding; construyo a diario con agentes manteniendo el criterio de ingeniería en el loop.

\vspace{8pt}
\begin{center}
    \textbf{Experiencia}
\end{center}
\vspace{-4pt}

\textbf{Proyecto Educa} \\
\textbf{Chief Technology Officer} \hfill \textit{Dic 2025 -- Presente}
\begin{itemize}
    \item \textbf{Dirijo el desarrollo de una plataforma de e-learning} para impartir y gestionar con escuelas los cursos de Proyecto Educa, dirigidos a estudiantes de 12 a 18 años.
    \item \textbf{Soy responsable de la arquitectura y la entrega técnica}, desde el modelo de datos y el backend hasta las prioridades de producto, los releases y el soporte en producción.
\end{itemize}

\vspace{10pt}
\textbf{Free2Z} \hfill \textit{Remoto} \\
\textbf{Principal Software Engineer} \hfill \textit{Abr 2025 -- Abr 2026}
\begin{itemize}
    \item \textbf{Envié features full-stack de punta a punta} integrando APIs en Node y Python/Django sobre un producto en vivo con usuarios reales.
    \item \textbf{Lideré la migración de un producto en producción de React a SvelteKit} sin romper el servicio, mejorando mantenibilidad y rendimiento.
\end{itemize}

\vspace{10pt}
\textbf{U.S. Department of State (via AVAA)} \\
\textbf{Software Engineer | Backend y Full-Stack} \hfill \textit{Nov 2024 -- Ago 2025}
\begin{itemize}
    \item \textbf{Único ingeniero de Billingua Talent}, dueño del arco completo sin project managers de por medio: requisitos, modelo de datos, backend, frontend, despliegue y cada release.
    \item \textbf{Construí una plataforma multi-tenant} con control de acceso por rol para cuatro tipos de usuario y aislamiento estricto por tenant, donde la integridad y la separación de los datos no eran negociables.
\end{itemize}

\vspace{10pt}
\textbf{Freelance} \hfill \textit{Remoto} \\
\textbf{Software Engineer} \hfill \textit{Sep 2024 -- Dic 2024}
\begin{itemize}
    \item \textbf{Entregué builds full-stack y scripts de automatización} en varios sectores, desde el descubrimiento y la priorización hasta el lanzamiento, la iteración y el soporte técnico.
\end{itemize}

\vspace{10pt}
\textbf{Asociación Venezolano Americana de Amistad (AVAA)} \\
\textbf{Software Developer} \hfill \textit{Oct 2022 -- Jul 2024}
\begin{itemize}
    \item \textbf{Propuse y construí SEP} (Node, PostgreSQL, Prisma sobre Azure), una plataforma multi-tenant que hoy sirve a alrededor de 300 becarios y staff en \textbf{3 capítulos nacionales}; empecé como pasante limpiando la data histórica, vi el problema real y me quedé como dueño técnico en producción.
    \item \textbf{Desarrollé SEA para EducationUSA}, centralizando historiales, servicios pagados, recibos, actividad mensual y reportes en un solo sistema operativo para los asesores.
\end{itemize}

\vspace{10pt}
\textbf{Asociación Venezolano Americana de Amistad (AVAA)} \hfill \textit{Caracas, Venezuela} \\
\textbf{Software Developer Intern} \hfill \textit{Sep 2021 -- Ago 2022}
\begin{itemize}
    \item \textbf{Construí un dashboard interactivo en Power BI} y escribí pipelines de datos en Python que mejoraron la confiabilidad del reporting y redujeron la limpieza manual.
\end{itemize}

\vspace{8pt}
\begin{center}
    \textbf{Voluntariado y liderazgo}
\end{center}
\vspace{-4pt}
\textbf{AVAA, Miembro del Comité de Educación} \hfill \textit{Sep 2024 -- Presente} \\
Administro SEP para cerca de 300 becarios en tres capítulos y apoyo mejoras de procesos, calidad de datos, reportes, rediseño curricular y propuestas institucionales. \\
\textbf{Dev3pack, Organizador del hub del Global Hackathon} \hfill \textit{Abr 2026 -- May 2026} \\
Organicé el hub de Caracas de un hackathon global de Web3 e IA, coordinando sede, agenda, aliados, mentores y participantes; también asesoré a los equipos en alcance y decisiones técnicas. \\
\textbf{JA Venezuela Alumni, Vicepresidente} \hfill \textit{Mar 2026 -- Presente} \\
Coordino iniciativas de formación, alianzas y comunidad con el equipo de alumni y organizaciones externas. \\
\textbf{Technovation Girls Venezuela, Mentor de programación} \hfill \textit{2026} \\
Guié a un equipo estudiantil en la definición del problema, el prototipado y el desarrollo de una aplicación móvil con MIT App Inventor.

\vspace{8pt}
\begin{center}
    \textbf{Educación}
\end{center}
\vspace{-4pt}
\textbf{Instituto de Estudios Superiores de Administración (IESA)} \hfill \textit{Caracas, Venezuela} \\
Maestría en Administración Pública (MPA) \hfill \textit{En curso} \\

\vspace{8pt}
\begin{center}
    \textbf{Idiomas}
\end{center}
\vspace{-4pt}
Español (nativo), Inglés (C1)

\end{document}
